\emph{(1).} The setting is a tame corank-one degeneration (\claimref{sheaf.tameness}). By clause (1) of \claimref{synthesis.wall-rung-count-is-height-of-dying-line-class} a three-rung completion gives a bundle $N$ of rank $2$ over $\Sig$ with $N(\overline G) \oplus \shL \oplus N \cong N(\overline G) \oplus \underline{\R}^{3}$ and $w(N) = (1 + x)^{-1}$ (Whitney, Fact \claimref{synthesis.imported-whitney-product-formula-and-projective-space}(i)), so $w_{1}(N) = x$; adding a complement of $N(\overline G)$ in a trivial bundle shows that $\shL \oplus N$ is stably trivial, and when $x^{2} = 0$, $(1 + x)^{-1} = 1 + x$. Choose a fibre metric on $N$; its unit circle bundle has fibre a circle, and the orientation system of $N$ is $\Z^{x}$ because $w_{1}(N) = x$. By (G+) a section over the one-skeleton exists, the obstruction to a section over the two-skeleton is a class in $H^{2}(\Sig; \Z^{x})$, and there is no further obstruction, the homotopy groups $\pi_{j}$ of the circle vanishing for $j \geq 2$. If $H^{2}(\Sig; \Z^{x}) = 0$, $N$ has a nowhere-zero section $s$ and $N = \R s \oplus N''$ with $N''$ a line bundle, $w_{1}(N'') = w_{1}(N) = x$. Then $N'' \otimes \shL$ has $w_{1} = 0$ and is trivial (Fact \claimref{sheaf.imported-real-line-bundles}), so $N'' \cong \shL$ and $N \cong \shL \oplus \underline{\R}$. Hence $N(\overline G) \oplus \shL \oplus \shL \oplus \underline{\R} \cong N(\overline G) \oplus \underline{\R}^{3}$, and $\shL \oplus \shL$ is stably trivial.

\emph{(2).} The universal cover $S^{2n+1}$ has the equivariant cell structure of Theorem \claimref{synthesis.cohomology-of-the-quarter-turn-lens-spaces-from-the-equivariant-cell-structure}(2): one free cell $e_{j}$ in each dimension, with boundary $(g - 1)e_{j-1}$ for $j$ odd and $(1 + g + g^{2} + g^{3})e_{j-1}$ for $j$ even. The line $\shL$ has holonomy $-1$ along the generator (in the stored walls it is spanned by $c(v)$ with $cJ = -Jc$), so equivariant cochains with values in $\Z^{x}$ are one copy of $\Z$ in each degree and the coboundary is the value of the boundary under $g \mapsto -1$: $-2$ from even to odd degrees and $1 - 1 + 1 - 1 = 0$ from odd to even degrees. So $H^{j}(\Sig_{n}; \Z^{x})$ is $\Z/2$ for $j$ odd, and for $j$ even it is the kernel of multiplication by $2$ (or, in no top even degree here, all of $\Z$ modulo $0$): zero, since the top degree $2n + 1$ is odd. With the group of order two the same computation gives $H^{2}(\R P^{2}; \Z^{x}) = \Z$ (the tangent bundle of the projective plane is a plane bundle with $w_{1} = x$ that is not $\shL \oplus \underline{\R}$) and $0$ for projective spaces of dimension at least three, as it must.

\emph{(3).} Let $n \geq 2$ and suppose $k(\Sig) \leq 3$; a completion with fewer rungs extends to one with three by adding standard hyperbolic planes, so a three-rung completion exists. By (2) and (1), $\shL \oplus \shL$ is stably trivial. For $n \geq 3$: the complexification of a stably trivial real bundle is stably trivial, and the complexification of $\shL \oplus \shL$ is twice the complexified dying line, $2\chi^{2}$ in the notation of clause (6) of the cited theorem, so $2(1 - \chi^{2}) = 0$ in $K(\Sig_{n}) = \Z[\chi]/(\chi^{4} - 1, (\chi - 1)^{n+1})$ (that clause, given the imported filtration, Fact \claimref{synthesis.imported-atiyah-hirzebruch-filtration-of-complex-k-theory}). With $s = \chi - 1$ the ideal of relations in $\Z[s]/(s^{n+1})$ has the triangular basis $s^{j}(4s + 6s^{2} + 4s^{3} + s^{4})$, $0 \leq j < n$; at $n = 3$, $2(1 - \chi^{2}) = -4s - 2s^{2}$ reduces to $-2s^{3}$, whose coefficient is not divisible by $4$, so it is not zero. For $n \geq 4$ the ring map $\chi \mapsto I$, $I = e_{1}e_{2}$ the unit bivector of a plane, to $\mathrm{span}_{\Z}\{1, I\}/((1 - I)^{n+1})$ is defined ($I^{4} = 1$) and sends $2(1 - \chi^{2})$ to $4 = -(1 - I)^{4}$, which is not a multiple of $(1 - I)^{5}$ because $4/(1 - I)^{5} = -(1 + I)/2$ is not integral; so $2(1 - \chi^{2}) \neq 0$. This contradicts stable triviality, so $k(\Sig) \geq 4$. For $n = 2$: by clause (3) of \claimref{synthesis.two-rungs-need-the-doubled-dying-line-stably-trivial}, given its hypothesis (L2), $q_{1}(\shL \oplus \shL) = 2y^{2} \neq 0$ in $H^{4}(\Sig_{2}; \Z) = \Z_{4}$ for both spin structures, and by its clause (2) $\shL \oplus \shL$ is then not stably trivial; so $k(\Sig) \geq 4$. The upper bounds are $e(\shL) \leq \dim\Sig + 1 = 6$ on $L^{5}(4)$ (clause (3) of the parent; $5$ with the explicit odd map of the companion result, which is not used as a premise) and $e(\shL) \leq 6$ on $L^{7}(4)$ (clause (3) of the two-rung result). The comparison: the bound $2^{k-1}(1 - \xi) = 0$ of \claimref{synthesis.lens-space-walls-need-more-rungs-than-the-mod-two-height} at $k = 3$ is $4(1 - \chi^{2}) = 0$, which holds for $n = 3, 4$, so that bound gives three there.

\emph{(4).} At $n = 1$, $H^{4}(\Sig_{1}) = 0$ and $2(1 - \chi^{2}) = 0$ in the K-ring, consistent with the stored two-rung completion of \claimref{synthesis.lens-space-wall-rung-count-sees-the-integral-twist}. With four rungs the complement has rank three; bundles of rank three with $w_{1} = x$ are not classified by a single class, and the argument does not extend.
